\section{Introduction}
\label{sec:intro}
Obtaining 3D representations from unposed sparse multi-view images in a feed-forward manner remains a fundamental challenge in computer vision and graphics, with broad implications for robotics~\cite{shen2023distilled,han2025d}, scene understanding~\cite{fedele2025superdec}, and novel view synthesis~\cite{ye2024no}. Recently, feed-forward 3D Gaussian splatting frameworks have gained considerable attention, demonstrating impressive performance in reconstruction and understanding~\cite{ye2024no,huang2025no,hong2024pf3plat,fan2024large,sheng2025spatialsplat}.

However, these approaches predominantly rely on \textit{dense, per-pixel} Gaussian predictions, which inherently leads to two critical issues: (1) it generates excessive redundant primitives that are often misaligned in 3D space~(Fig.~\ref{fig:concept}), and (2) it incurs substantial computational overhead when incorporating semantic features~\cite{sheng2025spatialsplat,fan2024large} through 2D-to-3D feature lifting. Consequently, prior works compress rich semantic information~\cite{simeoni2025dinov3,oquab2023dinov2,kim2025seg4diff} into lower-dimensional embeddings, resulting in information loss and sub-optimal scene understanding. This raises a fundamental question: \textbf{\textit{do we need such pixel-aligned Gaussians to reconstruct and understand 3D scenes?}}

As humans, we do not maintain pixel-perfect mental reconstructions of every surface to understand our surroundings. Instead, we form compact, semantically meaningful abstractions of identifying key objects, their rough spatial relationships, and overall scene structure~\cite{burgess2006spatial,shepard1971mental}. Drawing direct inspiration from human visual cognition, we propose a novel framework \oursname\ for learning \textbf{\textit{compact 3D Gaussians}} from unposed image observations in a feed-forward fashion. 

Similar to prior approaches~\cite{ye2024no,huang2025no}, we first extract image features from visual encoders with rich geometric priors (e.g., VGGT~\cite{wang2025vggt}). However, instead of learning to estimate per-pixel Gaussians directly from the extracted feature maps, we introduce a compact set of learnable query tokens that discover and decode essential 3D Gaussians. Specifically, we adopt a transformer architecture~\cite{vaswani2017attention} where learnable query tokens and the image features are processed through multiple self-attention blocks. We decode the refined learnable query tokens as 3D Gaussians, where the query tokens learn to aggregate essential information across multiple views to faithfully represent the scene.

Crucially, our framework requires no explicit supervision from ground-truth depths or scene decompositions. Despite training solely on photometric reconstruction, each token naturally learns to represent different regions, with each token attending to coherent spatial regions across views. This emergent behavior arises from the inherent structure of the task: to efficiently reconstruct novel views with a limited number of Gaussians, the model must learn to allocate Gaussians to meaningful regions. 

We show that after training, our model can estimate a compact set of 3D Gaussians, which enables efficient novel view synthesis while maintaining performance. In addition, the compact set of Gaussians enables 2D-to-3D feature lifting without compression, significantly improving 3D scene understanding tasks where the rich representation of the semantic features is critical.

\begin{figure*}[t]
    \centering
    \includegraphics[width=\linewidth]{figure/assets/figure2_concept.pdf}
    \vspace{-20pt}
    \caption{\textbf{Comparison of per-pixel and compact scene representations.} \textbf{(Left)}: Existing per-pixel estimators~\cite{ye2024no,huang2025no} predict one or multiple Gaussians per pixel, resulting in redundant Gaussians with misalignments across views. \textbf{(Right)}: Our method uses learnable Gaussian queries to discover and decode only compact 3D Gaussians at essential locations, achieving a compact representation with only 2K Gaussians and 4.1M memory while avoiding redundancy and achieving superior segmentation and novel view synthesis performance.}
    \label{fig:concept}
    \vspace{-10pt}
\end{figure*}



Our framework also provides a novel solution to a key challenge in 2D-to-3D feature lifting: handling multi-view feature inconsistencies. While other methods~\cite{zeid2025dino, lee2025cf3, cheng2024occam, marrie2025ludvig} require additional aggregation methods to handle inconsistent features across viewpoints, we observe that our model's emergent property of tokens attending to spatially coherent regions across views can be directly leveraged for feature aggregation. Specifically, we propose a view-invariant feature decoder~(\textbf{\oursF}) that reuses the attention maps from our learned Gaussian decoder~(\textbf{\ours}) while training only the value projections. This feature decoder can then take features from any visual encoder as input and decode multi-view aggregated features. By attaching these aggregated features to our estimated 3D Gaussians, we enable efficient novel view rendering with view-invariant features.

We validate the effectiveness of our approach through extensive experiments. For novel view synthesis, despite using \textbf{65$\times$} fewer Gaussians (only \textbf{2K}) than per-pixel methods~\cite{ye2024no,huang2025no}, we achieve competitive visual quality while enabling substantially faster rendering. More importantly, for 3D semantic understanding tasks, our compact Gaussians combined with the multi-view aggregated semantic features significantly outperform previous feed-forward approaches that attempt to lift view-inconsistent features through dense per-pixel Gaussians~\cite{fan2024large}. We further verify that our feature renderings can effectively replace previous feature upsamplers, while also outperforming previous view-invariant feature aggregation methods in multi-view correspondence evaluation.